% Generated draft: requires mathematical and source review.
\documentclass[11pt]{article}
\usepackage[margin=25mm]{geometry}
\usepackage{fontspec}
\setmainfont{Latin Modern Roman}
\usepackage{amsmath,amssymb,mathtools,unicode-math}
\setmathfont{Latin Modern Math}
\usepackage{enumitem,needspace,xurl,hyperref}
\setlength{\parindent}{0pt}
\setlength{\parskip}{5pt}
\setlength{\emergencystretch}{3em}
\newcommand{\N}{\mathbb{N}}
\newcommand{\Z}{\mathbb{Z}}
\newcommand{\Q}{\mathbb{Q}}
\newcommand{\R}{\mathbb{R}}
\newcommand{\C}{\mathbb{C}}
\newcommand{\F}{\mathbb{F}}
\newcommand{\A}{\mathbb{A}}
\newcommand{\PP}{\mathbb{P}}
\newcommand{\E}{\mathbb{E}}
\newcommand{\eps}{\varepsilon}
\newcommand{\dd}{\,\mathrm{d}}
\begin{document}
\section*{UnsolvedMath ID: 117}
\textbf{Source problem code:} catalogue entry
\subsection{Definitions and mathematical statement}
Let \(G\) be a finite abelian group, let \(d\geq 2\) be an integer, and let \(f:G\to\C\) satisfy \(|f(x)|\leq1\) for all \(x\in G\). Write \(\mathcal C z=\overline z\), and let expectations denote uniform averages. The Gowers norm is defined by
\[
\|f\|_{U^d(G)}^{2^d}=\E_{x,h_1,\ldots,h_d\in G}
\prod_{\omega\in\{0,1\}^d}\mathcal C^{|\omega|}
 f(x+\omega_1h_1+\cdots+\omega_dh_d).
\]
An inverse theorem turns a lower bound \(\|f\|_{U^d(G)}\geq\delta>0\) into correlation with a suitably specified structured object. If \(\mathcal S_d(G,B)\) denotes the chosen family of structured functions of complexity at most \(B\), its quantitative form would specify functions \(B_d\) and \(c_d\) such that
\[
\|f\|_{U^d(G)}\geq\delta
\ \Longrightarrow\
\exists\psi\in\mathcal S_d(G,B_d(\delta)):
\left|\E_{x\in G}f(x)\overline{\psi(x)}\right|\geq c_d(\delta).
\]
This is a framework for stating the question: the source does not define \(\mathcal S_d\), its complexity measure, the allowable group classes, or all permitted constant dependencies. Its request is to determine quantitative bounds, particularly a genuinely polynomial dependence on \(\delta^{-1}\) in unresolved general settings. Polynomial dependence would require the relevant complexity bounds to be at most \(C\delta^{-C}\) and the correlation bound at least \(C^{-1}\delta^C\), with \(C\) independent of the ambient size once the intended structural parameters are fixed. These schematic requirements do not supply the missing structural definitions.

The supplied background distinguishes finite-field problems for general bounded functions from inverse theorems over the integer interval \([N]=\{1,\ldots,N\}\), whose structured objects are nilsequences. It states that the integer nilsequence version has not been fully formulated at the desired polynomial strength. A finite-field formulation also requires the field, characteristic assumptions, and the intended class of polynomial phases, including any nonclassical phases, to be specified. No convention for the interval Gowers norm or nilsequence complexity is supplied here.

The supplied triage dated 2026-08-17 labels the agenda partially solved, although the registry status is open. It reports quasipolynomial \(U^{s+1}[N]\) inverse bounds of Leng--Sah--Sawhney for all \(s\geq3\); Mili\'cevi\'c's quasipolynomial \(U^4\) bounds in finite vector spaces and, in 2026, general finite abelian groups; and Milo--Moshkovitz's nearly polynomial results for polynomial inputs of degree \(d+1\) and related homogeneous inputs of degree below \(2d\). The last results concern restricted inputs and do not settle the general bounded-function problem. These are dated source assessments and do not constitute an independently verified current status.
\subsection{Short English statement}
Determine the quantitative dependence in inverse theorems for Gowers norms, aiming for polynomial bounds in the unresolved general cases. The supplied entry is an agenda spanning finite fields and integer intervals; a complete theorem requires a precise setting, structured comparison class, complexity measure, and dependence of constants.
\subsection{Sources}
\begin{itemize}
\item Local UnsolvedMath registry, permanent ID 117, source code catalogue entry, ``Inverse Theorem for Gowers Norms,'' including literature triage dated 2026-08-17. The source identifies this entry with Problem 56 of Ben Green's list by its statement, despite the local code catalogue entry.
\item B. Green, \emph{100 Open Problems}, Problem 56, supplied version retrieved 2026-08-17: \url{https://people.maths.ox.ac.uk/greenbj/papers/open-problems.pdf}.
\item J. Leng, A. Sah, M. Sawhney, \emph{Quasipolynomial bounds on the inverse theorem for the Gowers \(U^{s+1}[N]\)-norm}, arXiv:2402.17994 (2024): \url{https://arxiv.org/abs/2402.17994}.
\item L. Mili\'cevi\'c, \emph{Quasipolynomial inverse theorem for the \(U^4(\F_p^n)\) norm}, arXiv:2410.08966 (2024): \url{https://arxiv.org/abs/2410.08966}.
\item L. Mili\'cevi\'c, \emph{General inverse theory for the \(U^4\) norm}, arXiv:2601.01682 (2026): \url{https://arxiv.org/abs/2601.01682}.
\item T. Milo, G. Moshkovitz, \emph{Nearly-polynomial inverse theorem for the \(U^d\) norm in degree \(d+1\)}, arXiv:2603.16836v2 (2026): \url{https://arxiv.org/abs/2603.16836}.
\end{itemize}
\end{document}
